\documentclass[11pt]{article}
\usepackage[a4paper,margin=23mm,headheight=15pt]{geometry}
\usepackage[T1]{fontenc}
\usepackage[utf8]{inputenc}
\usepackage{lmodern,microtype}
\usepackage{amsmath,amssymb,mathtools,bm,mathrsfs}
\usepackage{booktabs,longtable,array,tabularx}
\usepackage{enumitem,xcolor,hyperref,fancyhdr,needspace}
\usepackage{xurl}
\definecolor{navy}{RGB}{25,54,88}
\definecolor{muted}{RGB}{72,89,108}
\definecolor{pale}{RGB}{237,243,249}
\definecolor{amber}{RGB}{252,245,226}
\hypersetup{colorlinks=true,linkcolor=navy,citecolor=navy,urlcolor=navy,
 pdftitle={From Blaha's Eq. 13 to a matrix-valued eDMFT stress formula},
 pdfauthor={Technical working note prepared for Kristjan Haule}}
\pagestyle{fancy}\fancyhf{}
\fancyhead[L]{\small\color{muted}LAPW stress: Eq. 13 to generalized G2}
\fancyhead[R]{\small\color{muted}6 November 2026}
\fancyfoot[C]{\thepage}
\renewcommand{\headrulewidth}{0.3pt}
\setlength{\parindent}{0pt}\setlength{\parskip}{5pt plus 1pt}
\setlist{nosep,leftmargin=1.6em}
\allowdisplaybreaks[2]\numberwithin{equation}{section}
\newcommand{\Tr}{\operatorname{Tr}}
\newcommand{\Rea}{\operatorname{Re}}
\newcommand{\Herm}{\operatorname{Herm}}
\newcommand{\diag}{\operatorname{diag}}
\newcommand{\dd}{\mathrm d}
\newcommand{\ii}{\mathrm i}
\newcommand{\kb}{\mathbf k}
\newcommand{\Id}{\mathbb I}
\newcommand{\e}{\epsilon}
\newcommand{\pa}{\partial_\mu}
\newcommand{\Es}{\mathcal E}
\newcommand{\Sig}{\widetilde\Sigma}
\newcommand{\rE}{{\rho_E}}
\newcommand{\bra}[1]{\langle #1|}
\newcommand{\ket}[1]{|#1\rangle}
\newcommand{\ip}[2]{\langle #1|#2\rangle}
\newcommand{\me}[3]{\langle #1|#2|#3\rangle}
\newcommand{\pathf}[1]{\nolinkurl{#1}}
\newcommand{\note}[2]{\par\medskip\noindent\colorbox{#1}{\parbox{\dimexpr\linewidth-2\fboxsep\relax}{#2}}\par\medskip}
\newcommand{\key}[1]{\note{pale}{#1}}
\newcommand{\caution}[1]{\note{amber}{#1}}
\newcommand{\newpart}[1]{\par\Needspace{12\baselineskip}\section{#1}}

\begin{document}
\hypersetup{pageanchor=false}
\begin{titlepage}
\centering
\vspace*{18mm}
{\Huge\bfseries\color{navy}Analytic Stress in\par DFT+eDMFT with LAPW\par}
\vspace{9mm}
{\LARGE From Blaha's Eq.\ (13), through Eq.\ (20),\par to a matrix-valued Appendix G2\par}
\vspace{16mm}
\begin{minipage}{0.88\linewidth}
This note derives the valence part of the stationary stress before using the
diagonal Kohn--Sham occupation identities. It retains both the density matrix
and its energy moment, makes the volume and representation cancellations
explicit, and gives six local matrix contractions corresponding to the six
lines of Blaha's Eq.\ (G2).

The normalized correlated projector is included in the functional derivative.
The last sections identify the source routines to re-audit, record the scope
of the available fixed-input tests, and distinguish derived identities from
unresolved implementation questions.

\textbf{Status: a derivation and audit reference.} This is not a claim that
the complete analytic eDMFT stress now agrees with the strain derivative of
the free energy. No program source was changed to prepare this document.
\end{minipage}
\vfill
{\large Prepared for Kristjan Haule\par}
\vspace{3mm}
{\large 6 November 2026\par}
\vspace{9mm}
{\small Primary DFT reference: Belbase, Tr\"oster, and Blaha,\par
\emph{Physical Review B} \textbf{104}, 174113 (2021).\par}
\end{titlepage}
\pagenumbering{roman}\hypersetup{pageanchor=true}
\tableofcontents
\clearpage\pagenumbering{arabic}

\newpart{The result and how to read it}
\label{sec:overview}
The starting point is the derivative of the stationary DFT+eDMFT functional,
not the derivative of the usual total-energy estimator. At the stationary
solution its valence, one-particle contribution is
\begin{equation}
 D_\mu^{\rm 1p}=\sum_\kb w_\kb\Rea\Tr_b
 \bigl[\rho_\kb\,\pa H_\kb-\rE_\kb\,\pa O_\kb\bigr].
 \label{eq:target}
\end{equation}
Here $\mu$ denotes a strain channel, $H$ is the KS Hamiltonian matrix,
$O$ is the LAPW overlap matrix in the chosen representation, and $\rho_E$
is the energy-weighted density matrix. Stress is $D_\mu/\Omega$.
The full derivative also contains the normalized-projector derivative,
the remaining density-functional and ion terms, and any explicitly
strain-dependent interaction parameters.

In a fixed central-band representation the valence contribution can be
organized as
\begin{equation}
 D_\mu^{\rm 1p}=\sum_\kb w_\kb\Rea\Tr_b
 \bigl[\rho(T^\mu+V^{\mu,\rm op}+B^\mu)
             -\rE O^{\mu,\eta}\bigr].
 \label{eq:overview}
\end{equation}
$T^\mu$ is the explicit kinetic-operator strain derivative;
$V^{\mu,\rm op}$ is the explicit potential-operator derivative in material
coordinates; $B^\mu$ and $O^{\mu,\eta}$ contain the volume-compensated
basis derivative. Their integral definitions are given in
Sec.~\ref{sec:matrices}. The superscript $\eta$ distinguishes this overlap
contribution from a derivative in an arbitrary moving representation.

\key{\textbf{The principal generalization of Eq.\ (20).}
The scalar DFT factor $f_n\epsilon_n$ must become the matrix $\rho_E$.
It cannot in general be replaced by $\tfrac12\{\Es,\rho\}$.
The full Hamiltonian and overlap bilinear forms must be retained before
contraction.}

Three points keep this statement from being mistaken for an ad hoc correction.
\begin{enumerate}
\item The normalization part of the first term in Blaha's Eq.\ (13)
      cancels the volume Jacobian. It is not a separate $\Tr M_\Sigma$
      pressure term.
\item A change of band representation contributes to the one-particle
      expression, but its pure representation part cancels against the
      corresponding projector change. A changing truncated subspace is
      a different matter.
\item The six expressions derived here describe the frozen-augmentation,
      nonrelativistic valence formula. A complete RELA implementation,
      radial-solution response, and the final total-stress assembly still
      require independent checks.
\end{enumerate}

Sections~\ref{sec:functional}--\ref{sec:general20} give the functional and
geometrical derivation. Sections~\ref{sec:local}--\ref{sec:g2} expand the result
into local radial and angular integrals. Sections~\ref{sec:projector} and
\ref{sec:audit} connect the mathematics to the next implementation audit.

\newpart{Notation, units, and what is held fixed}
\label{sec:notation}
\subsection{Indices and matrix spaces}
\begin{longtable}{>{$}p{0.19\linewidth}<{$}p{0.74\linewidth}}
\toprule
\kb,w_\kb & Brillouin-zone point and its fixed integration weight;
 $\partial_\mu w_\kb=0$. A spin block can be included in the index $\kb$;
 spin multiplicity is counted once, not again in the formulas.\\
m,n & Bands in the chosen valence space; this space may include the
 noninteracting bands outside the DMFT window.\\
K & LAPW basis index, including local-orbital basis functions when present.\\
a & Atom in the unit cell, with position $\tau_a$ and muffin-tin radius $R_a$.\\
L=(l,m_l) & Spherical-harmonic index. $m_l$ is a magnetic quantum number,
 distinct from a band index $m$.\\
i=(L,\lambda),j=(L',\lambda') & Local augmentation indices.
 For the elementary LAPW pair, $\lambda=0$ denotes $u_l$ and
 $\lambda=1$ denotes $\dot u_l=\partial u_l/\partial E$.\\
\ell & Correlated-orbital index, including the correlated site and spin
 structure as appropriate.\\
\mu=(\alpha\beta) & One of six symmetric strain channels;
 $\alpha,\beta\in\{x,y,z\}$.\\
\Omega & Unstrained cell volume.\\
\beta_T,T & Inverse temperature and temperature in energy units;
 $T=1/\beta_T$. We use $\beta_T$ to avoid confusion with a Cartesian index.\\
\omega_j,z_j & Fermionic frequency $\omega_j=(2j+1)\pi T$ and
 $z_j=\ii\omega_j+\mu_{\rm chem}$.\\
\bottomrule
\end{longtable}

The band matrices $H,O,G,\rho,\rho_E$ have size $N_b\times N_b$.
At the unstrained reference $O=\Id$ and
$H=\Es=\diag(\epsilon_1,\ldots,\epsilon_{N_b})$.
The local coefficient matrices $A,D^\mu$ have size $N_{\rm aug}\times N_b$.
The normalized projector $U$ and its raw counterpart $P$ have size
$N_b\times N_c$; the gradient $C$ has size $N_c\times N_b$.
Local matrices $S^x,H^x,V^x$ have size
$N_{\rm aug}\times N_{\rm aug}$.
$\Tr_b$ and $\Tr_c$ denote band and correlated-orbital traces.

We define $\Herm X=(X+X^\dagger)/2$, the anticommutator
$\{A,B\}=AB+BA$, and the commutator $[A,B]=AB-BA$.
The elementary trace convention is
$\Tr_b(\rho X)=\sum_{mn}\rho_{nm}X_{mn}$.
Unless shown explicitly, all matrices in a product have the same $\kb$
and spin block. No $k$ weight is absorbed into $\rho$ in this note.

\subsection{Strain and stress conventions}
Let $\boldsymbol\varepsilon$ be a symmetric strain tensor and define
\begin{equation}
 F=\Id+\boldsymbol\varepsilon,\qquad
 \mathbf r'=F\mathbf r,\qquad J=\det F,\qquad \Omega'=J\Omega.
\end{equation}
For a channel $\mu=(\alpha\beta)$ use the generator
\begin{equation}
 S^\mu_{ij}=\tfrac12(\delta_{i\alpha}\delta_{j\beta}
                         +\delta_{i\beta}\delta_{j\alpha}),\quad
 \partial_\mu F=S^\mu,\quad
 \gamma_\mu=\partial_\mu\ln J=\delta_{\alpha\beta}.
 \label{eq:strain}
\end{equation}
Thus the off-diagonal channel parameter is the engineering shear:
$F_{\alpha\beta}=F_{\beta\alpha}=s/2$. If a test instead sets both
off-diagonal entries equal to $s$, its derivative is twice the channel
derivative used here.

The thermodynamic sign convention is
\begin{equation}
 \sigma_\mu=\Omega^{-1}\partial_\mu\Gamma_{\rm stat},\qquad
 p=-\tfrac13\Tr\boldsymbol\sigma.
 \label{eq:stress-sign}
\end{equation}
For isotropic linear strain $F=(1+h)\Id$,
$\dd\Gamma/\dd h=\Omega(\sigma_{xx}+\sigma_{yy}+\sigma_{zz})$.
Divide by $3\Omega$ when comparing the hydrostatic mean stress.

Equations using $H=-\nabla^2/2+V$ are in \textbf{Hartree atomic units},
as in the paper. The code uses \textbf{Rydberg} energies: multiply
Hamiltonians, energy moments, and energy derivatives by two, while keeping
overlaps unchanged. In particular the inhomogeneous energy-derivative
radial equation must use the same energy variable as the radial function.
Approximately $1\,\mathrm{Ry}/a_0^3=14710.5\,\mathrm{GPa}$.

\subsection{Scope of the path}
The elementary derivation holds $R_a$ fixed in bohr, moves atomic centers
affinely, holds local axes fixed, and uses a fixed number of basis functions
and fixed band-index window. Radial functions are frozen as functions of
their physical radius and energy parameter, but their coordinate arguments
are strained. The distinction is crucial: frozen augmentation does
\emph{not} remove the $r u'(r)$ terms.

A finite-difference comparison must specify what is held fixed in the
potential, radial functions, projector, charge, self-energy, chemical
potential, and band window. A fixed-input matrix test and a relaxed
free-energy test are different derivatives.

\newpart{Stationary functional and the two density matrices}
\label{sec:functional}
\subsection{Which functional is differentiated?}
A useful auxiliary-field form of the DFT+eDMFT stationary functional is
\begin{align}
 \Gamma={}&\mathcal L[H_{\rm KS},O,U,\Sig]
 -\int_\Omega\dd^3r\,[V_{\rm KS}-V_{\rm ext}]\,n
 +E_{\rm Hxc}[n]+E_{\rm ion}\notag\\
 &+\Phi_{\rm imp}[G_{\rm loc}]-\Phi_{\rm DC}[n_{\rm loc}]
 -\Tr_{\omega,c}(\Sig G_{\rm loc})
 +\mu_{\rm chem}N .
 \label{eq:gamma}
\end{align}
The last term converts the grand potential to the fixed-$N$ Helmholtz
functional; omit it when working at fixed chemical potential.
$n(\mathbf r)$ is the charge density, $V_{\rm ext}$ the external ionic
potential, $E_{\rm Hxc}$ the Hartree plus exchange-correlation functional,
and $E_{\rm ion}$ the ionic electrostatic energy.
$\Phi_{\rm imp}$ is the local many-body functional and $\Phi_{\rm DC}$ its
double-counting subtraction. The local self-energy used for upfolding is
$\Sig=\Sigma_{\rm imp}-V_{\rm DC}$.
The trace $\Tr_{\omega,c}$ includes temperature and all Matsubara
frequencies as well as local orbital/site traces.

$\mathcal L$ is the \emph{regularized} lattice trace-log contribution
associated with
\begin{equation}
 G_\kb(z)^{-1}=zO_\kb-H_\kb-U_\kb\Sig(z)U_\kb^\dagger.
 \label{eq:dyson}
\end{equation}
It represents $-T\sum_{\kb j}w_\kb\ln\det[-G_\kb^{-1}(z_j)]$ with the
high-frequency reference and basis-measure normalization included.
The raw unnormalized determinant is not a definition of the thermodynamic
functional in a nonorthogonal basis.
Only its finite, consistently regularized derivative is needed below.

At stationarity, variations of the auxiliary fields and local Green
function cancel. In particular
$\delta(\Phi_{\rm imp}-\Phi_{\rm DC})
=\Tr_{\omega,c}(\Sig\delta G_{\rm loc})$
cancels the corresponding term in
$-\delta\Tr_{\omega,c}(\Sig G_{\rm loc})$;
the lattice trace-log supplies the opposite
$\Tr_{\omega,c}(G_{\rm loc}\delta\Sig)$.
The explicit KS-potential variation cancels through the density-functional
part. This statement does not discard changes of $U$ induced by a changing
KS potential, band subspace, or radial functions: these remain in the
specified total projector derivative. No impurity two-particle vertex
is required for this first derivative at stationarity.

\caution{Equation~\eqref{eq:gamma} is a stationary functional, not a claim
that every frozen-self-energy path is stationary. A numerical frozen-$\Sigma$
check must use the agreed linear impurity-functional correction and a
consistent charge-relaxation path.}

\subsection{Density and energy moment}
At the central orthonormal KS representation define
\begin{equation}
 \rho_\kb=T\sum_j^{\rm reg}G_\kb(z_j),\qquad
 \rE_\kb=T\sum_j^{\rm reg}z_jG_\kb(z_j).
 \label{eq:moments}
\end{equation}
The superscript ``reg'' includes the convergence prescription, the
high-frequency contact subtraction, and analytically summed reference.
It must be identical to that used in the trace log. It is fixed so that
for $\Sig=0$,
\begin{equation}
 \rho_{mn}=f_n\delta_{mn},\qquad
 (\rE)_{mn}=f_n\epsilon_n\delta_{mn},
 \label{eq:dftmoments}
\end{equation}
where $f_n$ is the Fermi occupation at the chosen temperature.
Both matrices in Eq.~\eqref{eq:moments} are Hermitian after the complete
frequency sum.

The finite derivative of the lattice term is
\begin{equation}
 \delta\mathcal L=\sum_\kb w_\kb\Rea\Tr_b
 \left(\rho\delta H-\rE\delta O
       +T\sum_j^{\rm reg}G(z_j)\,\delta\Sigma^{\rm emb}(z_j)\right),
 \label{eq:logderiv}
\end{equation}
with $\Sigma^{\rm emb}=U\Sig U^\dagger$. The part involving
$\delta\Sig$ cancels in Eq.~\eqref{eq:gamma}; the explicit $\delta U$
part remains. This gives Eq.~\eqref{eq:target} and the projector term
derived in Sec.~\ref{sec:projector}.

\subsection{What is \texorpdfstring{$M_\Sigma$}{M-Sigma}?}
Left and right Dyson identities give
\begin{align}
 B_L&=\rE-\Es\rho
     =T\sum_j^{\rm reg}\Sigma^{\rm emb}(z_j)G(z_j),\\
 B_R&=\rE-\rho\Es=B_L^\dagger,\\
 M_\Sigma&=\tfrac12(B_L+B_R)
   =\rE-\tfrac12\{\Es,\rho\}.
 \label{eq:msigma}
\end{align}
The symbol $B_L$ is a Dyson residual moment, not the basis derivative
matrix $B^\mu$ introduced later. Explicitly,
\begin{equation}
 M_\Sigma=\frac T2\sum_j^{\rm reg}
       [\Sigma^{\rm emb}G+G\Sigma^{\rm emb}].
\end{equation}
$M_\Sigma$ has energy units and vanishes in the noninteracting limit
with the same Hamiltonian reference. The decomposition
\begin{equation}
 \rE=\tfrac12\{\Es,\rho\}+M_\Sigma
 \label{eq:momentdecomp}
\end{equation}
is exact. Replacing $\rho_E$ by a symmetrized band energy omits
$M_\Sigma$ and is not justified by merely Hermitizing the DFT formula.

\newpart{Recovering Blaha's Eq. (13) without discarding a term}
\label{sec:eq13}
Let $\phi_K[\varepsilon](\mathbf r')$ be a strained LAPW basis function and
$c_{Kn}[\varepsilon]$ a coefficient. The material-coordinate wave function is
\begin{equation}
 \widetilde\psi_n[\varepsilon](\mathbf r)
 =\psi_n[\varepsilon](F\mathbf r)
 =\sum_K c_{Kn}[\varepsilon]\,
             \phi_K[\varepsilon](F\mathbf r).
\end{equation}
All derivatives below are evaluated at zero strain; a tilde is suppressed
after this definition. Define three derivatives:
\begin{align}
 b_n^\mu(\mathbf r)&=\sum_K(\pa c_{Kn})\phi_K(\mathbf r),
 \label{eq:bdef}\\
 \eta_n^\mu(\mathbf r)&=\sum_K\frac{c_{Kn}}{\sqrt\Omega}
 \pa\!\left[\sqrt{\Omega[\varepsilon]}
           \phi_K[\varepsilon](F\mathbf r)\right]_0,
 \label{eq:etadef}\\
 \zeta_n^\mu&=b_n^\mu-\frac{\gamma_\mu}{2}\psi_n.
\end{align}
The product rule now gives exactly the separation in the paper:
\begin{equation}
 \boxed{\pa\psi_n=b_n^\mu-\frac{\gamma_\mu}{2}\psi_n+\eta_n^\mu
                   =\zeta_n^\mu+\eta_n^\mu.}
 \label{eq:13}
\end{equation}
Blaha's first term is $\zeta_n^\mu$, not just $b_n^\mu$.
It combines the derivative of the expansion coefficients and the derivative
of $\Omega^{-1/2}$. His second term is $\eta_n^\mu$.

\subsection{Why the first term vanishes in the DFT derivation}
In a finite LAPW basis the KS eigenproblem implies
\begin{equation}
 \me{\phi_K}{H_{\rm KS}-\epsilon_n}{\psi_n}=0\quad\hbox{for every }K.
 \label{eq:galerkin}
\end{equation}
This is a projected variational equation; it does not say that the
finite-basis wave function solves the continuum differential equation
pointwise. Since $\zeta_n^\mu$ lies in the same basis space,
$\me{\zeta_n^\mu}{H-\epsilon_n}{\psi_n}=0$.
That is the precise reason for the sentence below Eq.\ (13) in the paper.

In DMFT the relevant contraction contains $\rho$ and $\rho_E$, not a
separate factor $f_n(H-\epsilon_n)$ for each state. One must therefore
return to the matrix derivative before making this simplification.
The next section does so, including the Jacobian that is already canceled
in the DFT Rayleigh quotient leading to the paper's Eq.\ (12).

\subsection{Interstitial behavior}
For an interstitial plane wave,
\begin{equation}
 \sqrt{\Omega[\varepsilon]}\phi_{\kb K}[\varepsilon](F\mathbf r)
 =\exp[\ii F^{-T}(\kb+\mathbf K)\cdot F\mathbf r]
 =\exp[\ii(\kb+\mathbf K)\cdot\mathbf r].
 \label{eq:plane}
\end{equation}
Hence $\eta_n^\mu=0$ there. This is the content of Appendix D1.
It does \emph{not} set the interstitial kinetic stress to zero: the
kinetic operator itself changes under the pullback.

\newpart{The volume cancellation and generalized Eq. (20)}
\label{sec:general20}
\subsection{Differentiate the matrix integrals first}
Write $\dot\psi_n=\pa\psi_n$ and let
$H^{\mu,\rm op}$ denote the derivative of the pulled-back differential
operator, with the wave functions held fixed. Differentiating the
whole-cell matrix integrals gives
\begin{align}
 (\pa H)_{mn}={}&\gamma_\mu H_{mn}
 +\me{\dot\psi_m}{H}{\psi_n}
 +\me{\psi_m}{H}{\dot\psi_n}
 +\me{\psi_m}{H^{\mu,\rm op}}{\psi_n},
 \label{eq:Hprime}\\
 (\pa O)_{mn}={}&\gamma_\mu O_{mn}
 +\ip{\dot\psi_m}{\psi_n}+\ip{\psi_m}{\dot\psi_n}.
 \label{eq:Oprime}
\end{align}
These formulas use the complete periodic cell. If a calculation instead
differentiates separate moving MT/interstitial domains, their boundary
terms must reconstruct these identities. The elementary formulas assume
continuous LAPW functions with the required matching; APW derivative jumps
and associated surface terms need separate care.

For a set of variations $x_n$ define the real linear contraction
\begin{equation}
 \mathcal P_\mu[x]=2\Rea\sum_{mn}\left[
 \rho_{nm}\me{x_m}{H}{\psi_n}
 -(\rE)_{nm}\ip{x_m}{\psi_n}\right].
 \label{eq:Pdef}
\end{equation}
At one $\kb$, Eq.~\eqref{eq:target} becomes
\begin{equation}
 D_\mu^{\rm 1p}=\gamma_\mu\Tr_b(\rho H-\rE O)
                +\mathcal P_\mu[\dot\psi]
                +\Rea\Tr_b(\rho H^{\mu,\rm op}).
 \label{eq:prevolume}
\end{equation}
Hermiticity makes the displayed traces real where appropriate.

Insert Eq.~\eqref{eq:13}. Its normalization contribution satisfies
\begin{equation}
 \mathcal P_\mu[-\gamma_\mu\psi/2]
       =-\gamma_\mu\Tr_b(\rho H-\rE O).
 \label{eq:volcancel}
\end{equation}
Therefore
\begin{equation}
 \boxed{D_\mu^{\rm 1p}=\Rea\Tr_b(\rho H^{\mu,\rm op})
                        +\mathcal P_\mu[\eta]
                        +\mathcal P_\mu[b].}
 \label{eq:postvolume}
\end{equation}
The cancellation is valid also when $[\rho,\Es]\ne0$ and
$M_\Sigma\ne0$. At the central orthonormal reference,
$\Tr(\rho H-\rE)=-\Tr M_\Sigma$. Restoring only the
$-\gamma\psi/2$ piece would produce $+\gamma\Tr M_\Sigma$;
the omitted Jacobian is exactly its negative.

\key{\textbf{No standalone volume correction follows from Eq.\ (13).}
Keep the Jacobian and wave-function normalization together. Their
cancellation is a whole-cell identity and must precede restricting
$\eta$ to muffin-tin integrals.}

\subsection{The matrix-valued valence correction}
With fixed central coefficients, $b=0$. Since $\eta=0$ in the
interstitial, the generalized counterpart of the paper's Eq.\ (20) is
\begin{equation}
 \boxed{\sigma_\mu^{\rm corr}=
 \frac{2}{\Omega}\Rea\sum_{a\kb}w_\kb\sum_{mn}
 \left[\rho_{nm}\me{\eta_m^\mu}{H_{\rm KS}}{\psi_n}_a
       -(\rE)_{nm}\ip{\eta_m^\mu}{\psi_n}_a\right].}
 \label{eq:20dmft}
\end{equation}
The subscript $a$ means integration over the central MT sphere of atom $a$.
For Eq.~\eqref{eq:dftmoments}, this reduces to
\begin{equation}
 \sigma_\mu^{\rm corr,DFT}
 =\frac{2}{\Omega}\Rea\sum_{a\kb n}w_\kb f_n
       \me{\eta_n^\mu}{H_{\rm KS}-\epsilon_n}{\psi_n}_a,
 \label{eq:20dft}
\end{equation}
which is precisely the structure of Eq.\ (20).

The potential in the matrix element in Eq.~\eqref{eq:20dmft} is part
of the basis derivative and must remain there. Cancellation of the
\emph{induced potential variation} in the stationary functional is not
permission to delete this potential from the Pulay matrix element.

\newpart{Band representation versus physical subspace response}
\label{sec:representation}
Suppose a part of the coefficient response can be expressed inside the
chosen band span as
\begin{equation}
 b_n=\sum_m\psi_m L_{mn}.
\end{equation}
Then the term retained in Eq.~\eqref{eq:postvolume} is
\begin{align}
 \mathcal P_\mu[b]
 &=2\Rea\Tr_b[L^\dagger(\Es\rho-\rE)]\notag\\
 &=2\Rea\Tr_b\!\left[L^\dagger
           \left(\tfrac12[\Es,\rho]-M_\Sigma\right)\right]
 =-2\Rea\Tr_b(L^\dagger B_L).
 \label{eq:bterm}
\end{align}
This vanishes in DFT. It generally does not vanish in the DMFT
one-particle term alone. That fact does not establish a new physical
stress, because the same change of representation also changes $U$.

Let
\begin{equation}
 C=T\sum_j^{\rm reg}\Sig(z_j)U^\dagger G(z_j),\qquad
 UC=B_L.
 \label{eq:C}
\end{equation}
For a pure orthonormal change of band frame, $L^\dagger=-L$ and
$\delta U=L^\dagger U$. Its projector contribution is
\begin{equation}
 \delta D^{\rm proj}=2\Rea\Tr_c(C\delta U)
      =2\Rea\Tr_b(L^\dagger UC)
      =2\Rea\Tr_b(L^\dagger B_L),
 \label{eq:gaugecancel}
\end{equation}
which cancels Eq.~\eqref{eq:bterm}. The same covariance holds for a
nonorthogonal frame when its overlap metric and projector normalization
are transformed consistently; see Sec.~\ref{sec:projector}.

\subsection{Two consistent ways to organize the calculation}
\begin{enumerate}
\item \textbf{Fixed reference coefficients.} Use the central coefficients
      to build $\partial H$ and $\partial O$ as basis derivatives.
      Then $b=0$ in these matrix derivatives. Compute the projector
      derivative in the same representation, including metric factors.
\item \textbf{A moving KS basis.} Retain its coefficient response in both
      the one-particle and projector derivatives. Remove the common
      pure-frame contribution only after demonstrating the cancellation.
\end{enumerate}
Mixing the first convention for one part and the second for another
can leave a large, unphysical residual even when the separate matrix
derivatives are accurate.

For a fixed-index window $W_\kb$, rotation entirely within $W_\kb$ is a
change of frame. Mixing with states outside $W_\kb$ changes the physical
subspace selected by the window and cannot be discarded by the preceding
argument. A constant number of bands does not imply
$\partial_\mu\Pi_{W_\kb}=0$, where $\Pi_{W_\kb}$ is its projector.
If the window cuts a degeneracy, this derivative may not be uniquely
smooth; include the degenerate multiplet or otherwise define a smooth
subspace.

\newpart{Local form of the basis derivative}
\label{sec:local}
\subsection{Augmentation coefficients and radial functions}
Within atom $a$ use $\mathbf r_a=\mathbf r-\tau_a$,
$r=|\mathbf r_a|$, and $\widehat{\mathbf r}=\mathbf r_a/r$. Write
\begin{equation}
 f_i^a(\mathbf r_a)=u_i^a(r)Y_{L_i}(\widehat{\mathbf r}),\qquad
 \psi_n(\mathbf r_a)=\sum_i A^a_{in}f_i^a(\mathbf r_a),\qquad
 A^a_{in}=\sum_K c_{Kn}a^a_{iK}.
 \label{eq:localA}
\end{equation}
$a_{iK}$ are the LAPW matching coefficients. The symbol $A$ here is
Blaha's local augmentation coefficient, not the full LAPW eigenvector $c$.
For the basic LAPW pair
\begin{equation}
 u_{l0}=u_l(r,E_l),\qquad u_{l1}=\dot u_l(r,E_l),
 \qquad \dot u_l=\partial u_l/\partial E.
\end{equation}
A prime always means $\dd/\dd r$; a dot means an energy derivative.
The spherical-harmonic normalization is
$\int\dd\widehat{\mathbf r}\,Y_L^*Y_{L'}=\delta_{LL'}$.

Define the volume-compensated matching derivative at fixed $c$:
\begin{equation}
 D^{a\mu}_{in}=\sum_K\frac{c_{Kn}}{\sqrt\Omega}
        \pa[\sqrt{\Omega[\varepsilon]}a^a_{iK}[\varepsilon]]_0.
 \label{eq:D}
\end{equation}
This is Blaha's G3--G4. If $s_{in}^\mu$ is the \emph{physical}
fixed-$c$ derivative of $A_{in}$, including $\Omega^{-1/2}$, then
\begin{equation}
 D^\mu=s^\mu+\tfrac12\gamma_\mu A.
 \label{eq:DvsS}
\end{equation}
One must not identify a routine returning $s^\mu$ with a routine
returning $D^\mu$ without this conversion.

\subsection{Radial and angular coordinate derivatives}
Define the scalar angular multiplier and angular differential operator
\begin{equation}
 p_\mu(\widehat{\mathbf r})=\hat r_\alpha\hat r_\beta,\qquad
 \mathscr A_\mu=\tfrac12
  (\hat r_\beta\nabla_{\Omega,\alpha}
       +\hat r_\alpha\nabla_{\Omega,\beta}).
 \label{eq:angular}
\end{equation}
$\boldsymbol\nabla_\Omega$ is the dimensionless tangential gradient
on the unit sphere. Since
$\boldsymbol\nabla=\widehat{\mathbf r}\,\partial_r
                     +r^{-1}\boldsymbol\nabla_\Omega$,
\begin{equation}
 (S^\mu\mathbf r_a)\cdot\boldsymbol\nabla
          [u_i(r)Y_i(\widehat{\mathbf r})]
 =r u_i'(r)p_\mu Y_i+u_i(r)\mathscr A_\mu Y_i.
 \label{eq:coordderiv}
\end{equation}
Here $Y_i$ abbreviates $Y_{L_i}$.
For isotropic dilation, summing the three diagonal channels gives
$\sum_\alpha p_{\alpha\alpha}=1$ and
$\sum_\alpha\mathscr A_{\alpha\alpha}=0$.

Combining the matching and coordinate derivatives gives
\begin{equation}
 \boxed{\eta_n^\mu=\sum_i\left[
 D^\mu_{in}u_iY_i
 +A_{in}\,r u_i'p_\mu Y_i
 +A_{in}\,u_i\mathscr A_\mu Y_i\right].}
 \label{eq:eta3}
\end{equation}
These are the three channels that produce the six terms in G2 when
$H_{\rm KS}=H_{\rm sph}+V_{\rm ns}$ is used.

\subsection{What is omitted by frozen augmentation?}
An explicit response of the radial solution would add
\begin{equation}
 \eta_n^{\mu,\rm rad\,resp}
       =\sum_i A_{in}(\pa u_i)_{r\,\rm fixed}Y_i,
 \label{eq:radialresponse}
\end{equation}
together with the consistent response of matching coefficients,
linearization energies, local orbitals, and normalized projector.
This is distinct from the $r u_i'$ term already in Eq.~\eqref{eq:eta3}.
The present note does not establish that this omitted response is the
source of the remaining total-stress discrepancy.

\newpart{Local integral matrices: definitions without diagonal assumptions}
\label{sec:integrals}
Suppress $a,\mu$ temporarily and introduce three bra functions:
\begin{equation}
 g_i^0=u_iY_i,\qquad
 g_i^r=r u_i'p_\mu Y_i,\qquad
 g_i^{\rm ang}=u_i\mathscr A_\mu Y_i.
 \label{eq:gfunctions}
\end{equation}
For $x\in\{0,r,\mathrm{ang}\}$ define
\begin{align}
 S^x_{ij}&=\int_{r<R_a}\dd^3r\,(g_i^x)^*f_j,\label{eq:Slocal}\\
 H^x_{ij}&=\int_{r<R_a}\dd^3r\,(g_i^x)^*H_{\rm sph}f_j,
              \label{eq:Hlocal}\\
 V^x_{ij}&=\int_{r<R_a}\dd^3r\,(g_i^x)^*V_{\rm ns}f_j.
              \label{eq:Vlocal}
\end{align}
These are \emph{one-sided} bilinear forms, with the differential
Hamiltonian acting on the ket. They need not individually be Hermitian.
The final band matrices are Hermitian sums, not individually symmetrized
local factors. Integration by parts changes the radial volume matrix
and introduces a surface term; both must be changed together.

\subsection{Overlap integrals}
The three local overlaps are
\begin{align}
 S^0_{ij}&=\delta_{L_iL_j}\int_0^{R_a}r^2\dd r\,u_i^*u_j,
                  \label{eq:S0}\\
 S^r_{ij}&=\left[\int_0^{R_a}r^3\dd r\,u_i'^*u_j\right]
              \int\dd\widehat{\mathbf r}\,Y_i^*p_\mu Y_j,
                  \label{eq:Sr}\\
 S^{\rm ang}_{ij}&=\left[\int_0^{R_a}r^2\dd r\,u_i^*u_j\right]
              \int\dd\widehat{\mathbf r}\,
                          (\mathscr A_\mu Y_i)^*Y_j.
                  \label{eq:Sa}
\end{align}
The angular derivative can connect different $l$ values. The radial
integral in Eq.~\eqref{eq:Sa} must therefore be evaluated for its actual
pair $(l_i,l_j)$. The familiar same-$l$ identities
$\langle u_l|u_l\rangle=1$ and
$\langle u_l|\dot u_l\rangle=0$ cannot be applied across different $l$.

\subsection{Spherical Hamiltonian action}
In Hartree units the radial equations imply
\begin{equation}
 H_{
m sph}(u_lY_L)=E_lu_lY_L,\qquad
 H_{
m sph}(\dot u_lY_L)=E_l\dot u_lY_L+u_lY_L.
 \label{eq:radialeq}
\end{equation}
For any channel $x$,
\begin{equation}
 \boxed{H^x_{ij}=E_{l_j}S^x_{ij}
                  +\delta_{\lambda_j,1}S^x_{i,j_0},}
 \label{eq:HfromS}
\end{equation}
where $j_0=(L_j,0)$. No band eigenvalue has entered this local identity.
It is therefore safe to use before taking an interacting density-matrix
trace. Actual radial-normalization and energy-unit conventions must
match Eq.~\eqref{eq:radialeq}.

\subsection{Nonspherical potential integrals}
Write
\begin{equation}
 V_{\rm ns}(\mathbf r)=\sum_{L''\ne0}V_{L''}(r)Y_{L''}(\widehat{\mathbf r}).
\end{equation}
Then
\begin{align}
 V^0_{ij}&=\sum_{L''\ne0}
 \left[\int_0^{R_a}r^2\dd r\,u_i^*V_{L''}u_j\right]
 \int\dd\widehat{\mathbf r}\,Y_i^*Y_{L''}Y_j,\label{eq:V0}\\
 V^r_{ij}&=\sum_{L''\ne0}
 \left[\int_0^{R_a}r^3\dd r\,u_i'^*V_{L''}u_j\right]
 \int\dd\widehat{\mathbf r}\,Y_i^*p_\mu Y_{L''}Y_j,\label{eq:Vr}\\
 V^{\rm ang}_{ij}&=\sum_{L''\ne0}
 \left[\int_0^{R_a}r^2\dd r\,u_i^*V_{L''}u_j\right]
 \int\dd\widehat{\mathbf r}\,(\mathscr A_\mu Y_i)^*Y_{L''}Y_j.
 \label{eq:Va}
\end{align}
Expanding $p_\mu$ and $\mathscr A_\mu Y_i$ in spherical harmonics gives
the Gaunt-coefficient sums of G2. The integral form fixes all conjugations
and derivative placements without depending on a particular Gaunt table.

\newpart{The six generalized lines of Appendix G2}
\label{sec:g2}
Insert Eq.~\eqref{eq:eta3} into Eq.~\eqref{eq:20dmft}, split
$H=H_{\rm sph}+V_{\rm ns}$, and use the local matrices above.
For one atom and one $\kb$, define the six real energy derivatives:
\begin{align}
 d_1^\mu&=2\Rea\Tr_b[\rho D^{\mu\dagger}H^0A
                             -\rE D^{\mu\dagger}S^0A],\label{eq:g21}\\
 d_2^\mu&=2\Rea\Tr_b[\rho A^\dagger H^rA
                             -\rE A^\dagger S^rA],\label{eq:g22}\\
 d_3^\mu&=2\Rea\Tr_b[\rho A^\dagger H^{\rm ang}A
                             -\rE A^\dagger S^{\rm ang}A],\label{eq:g23}\\
 d_4^\mu&=2\Rea\Tr_b[\rho D^{\mu\dagger}V^0A],\label{eq:g24}\\
 d_5^\mu&=2\Rea\Tr_b[\rho A^\dagger V^rA],\label{eq:g25}\\
 d_6^\mu&=2\Rea\Tr_b[\rho A^\dagger V^{\rm ang}A].\label{eq:g26}
\end{align}
Their sum is
\begin{equation}
 \boxed{\sigma_\mu^{\rm corr}
 =\frac1\Omega\sum_{a\kb}w_\kb\sum_{q=1}^6d_q^{a\kb,\mu}.}
 \label{eq:G2full}
\end{equation}
The $H^x,S^x,V^x$ for $x=r,\mathrm{ang}$ depend on $\mu$ through
their angular functions even though this superscript is suppressed.

\begin{tabularx}{\linewidth}{@{}c l X@{}}
\toprule
G2 line & Basis derivative & Local operator / overlap content\\
\midrule
1 & Matching, $D^\mu$ & $H^0$ and $S^0$\\
2 & Radial coordinate, $r u'$ & $H^r$ and $S^r$\\
3 & Angular coordinate & $H^{\rm ang}$ and $S^{\rm ang}$\\
4 & Matching, $D^\mu$ & Nonspherical potential $V^0$\\
5 & Radial coordinate, $r u'$ & Nonspherical potential $V^r$\\
6 & Angular coordinate & Nonspherical potential $V^{\rm ang}$\\
\bottomrule
\end{tabularx}

\subsection{An explicit index expression}
The first line, after using the radial equation, is
\begin{align}
 d_1^\mu=2\Rea\sum_{mnij}(D^\mu_{im})^*A_{jn}
 \Big[&(E_{l_j}\rho_{nm}-(\rE)_{nm})S^0_{ij}\notag\\
       &+\rho_{nm}\delta_{\lambda_j,1}S^0_{i,j_0}\Big].
 \label{eq:g21index}
\end{align}
The second and third lines have exactly the same bracket, with
$D^\mu_{im}$ replaced by $A_{im}$ and $S^0$ replaced by $S^r$ or
$S^{\rm ang}$. In DFT, $\rho_{nm}=f_n\delta_{nm}$ and
$(\rE)_{nm}=f_n\epsilon_n\delta_{nm}$, so the bracket contains
$f_n(E_{l_j}-\epsilon_n)$ and recovers the corresponding G2 factors.

The off-diagonal generalization therefore does require changing one
band index on the bra relative to the ket, as suggested in the discussion.
It also requires replacing the entire energy-weighted occupation by
$\rho_E$, including $M_\Sigma$. The latter information cannot be
recovered from occupations alone.

\subsection{Where the symmetric combination belongs}
For any Hermitian $\rho$ and any band matrix $X$,
\begin{equation}
 2\Rea\Tr_b(\rho X)=\Tr_b[\rho(X+X^\dagger)]
                         =2\Tr_b(\rho\Herm X).
 \label{eq:hermitize}
\end{equation}
For example, the Hermitian matching-potential band matrix is
\begin{equation}
 D^\dagger V^0A+A^\dagger(V^0)^\dagger D.
\end{equation}
In components this produces the paired $D_m^*A_n$ and $A_m^*D_n$
terms. Using $\Herm X=(X+X^\dagger)/2$ is equally correct if the
external factor two is retained. Applying a half and then dropping that
factor loses half the contribution. Hermitize the \emph{complete band
matrix}, not a local radial factor whose bra and ket play different roles.

\newpart{Explicit band matrices \texorpdfstring{$T^\mu$, $B^\mu$, and $O^\mu$}{T, B, and O}}
\label{sec:matrices}
Define the one-sided band matrices
\begin{align}
 Y^\mu={}&\sum_a\left[D^{a\mu\dagger}S^{a0}A^a
       +A^{a\dagger}S^{ar}A^a+A^{a\dagger}S^{a,\rm ang}A^a\right],\\
 Z^\mu={}&\sum_a\left[D^{a\mu\dagger}(H^{a0}+V^{a0})A^a
       +A^{a\dagger}(H^{ar}+V^{ar})A^a\right.\notag\\
       &\hspace{29mm}\left.+A^{a\dagger}(H^{a,\rm ang}
                                            +V^{a,\rm ang})A^a\right].
\end{align}
Then
\begin{equation}
 \boxed{O^{\mu,\eta}=Y^\mu+Y^{\mu\dagger},\qquad
        B^\mu=Z^\mu+Z^{\mu\dagger}.}
 \label{eq:BO}
\end{equation}
Equation~\eqref{eq:G2full} is equivalently
\begin{equation}
 \sigma_\mu^{\rm corr}=\frac1\Omega\sum_\kb w_\kb
 \Rea\Tr_b[\rho B^\mu-\rE O^{\mu,\eta}].
\end{equation}
These are definitions of the full off-diagonal matrices, not just their
diagonal entries. They identify which local information is needed even
if an implementation contracts the matrices without storing them.

For the explicit kinetic operator,
\begin{equation}
 T^\mu_{mn}=\frac12\int_\Omega\dd^3r\,
 \psi_m^*(\partial_\alpha\partial_\beta
                  +\partial_\beta\partial_\alpha)\psi_n.
 \label{eq:T}
\end{equation}
This is the off-diagonal bilinear form underlying Eqs.\ (19) and (G1).
For smooth periodic functions, integration by parts gives the equivalent
form
\begin{equation}
 T^\mu_{mn}=-\frac12\int_\Omega\dd^3r\,
 [(\partial_\alpha\psi_m)^*(\partial_\beta\psi_n)
  +(\partial_\beta\psi_m)^*(\partial_\alpha\psi_n)].
 \label{eq:Tgrad}
\end{equation}
When splitting the domain or using derivative-discontinuous basis
functions, keep the corresponding surface contributions consistently.

The explicit potential-operator matrix is
\begin{equation}
 V^{\mu,\rm op}_{mn}=
 \int_\Omega\dd^3r\,\psi_m^*(\mathbf r)
 \pa[V_{\rm KS}[\varepsilon](F\mathbf r)]_0\psi_n(\mathbf r).
\end{equation}
In the fixed-$c$, whole-cell convention, within the specified basis
approximation,
\begin{equation}
 \boxed{\pa O=O^{\mu,\eta},\qquad
        \pa H=T^\mu+V^{\mu,\rm op}+B^\mu.}
 \label{eq:matrixderivs}
\end{equation}
A moving representation adds $L^\dagger+L$ to $\partial O$ and
$L^\dagger\Es+\Es L$ to $\partial H$. These additions lead to
Eq.~\eqref{eq:bterm}.

\subsection{What an energy-linear ``kernel'' means}
A diagonal DFT routine can encode a bilinear form as
\begin{equation}
 K^\mu(E)=T^\mu+B^\mu-E O^{\mu,\eta}.
 \label{eq:kernel}
\end{equation}
Here ``kernel'' means a matrix representing this quadratic/bilinear form;
it is not a two-particle vertex. If Eq.~\eqref{eq:kernel} has been
established for the \emph{full} matrix, then in a chosen energy unit
$K^\mu(0)-K^\mu(1)=O^{\mu,\eta}$.
Knowing only $\langle n|K^\mu(\epsilon_n)|n\rangle$ does not prove
that this subtraction reconstructs the off-diagonal overlap derivative.
Rearrangements valid on KS eigenstates may remove information needed
for that inference. The integral definitions above provide the direct
audit target.

\newpart{Normalized projectors and consistency of the representation}
\label{sec:projector}
\subsection{The projector is not just a raw overlap}
For raw projected local orbitals $\varphi_\ell$, define in an
orthonormal window basis
\begin{equation}
 P_{n\ell}(\kb)=\ip{\psi_{n\kb}}{\varphi_\ell},\quad
 \mathcal S=\sum_\kb w_\kb P_\kb^\dagger P_\kb,\quad
 X=\mathcal S^{-1/2},\quad U_\kb=P_\kb X.
 \label{eq:U}
\end{equation}
This notation represents one normalization block; site/symmetry blocks
must be treated as in the actual projector construction.
The metric $\mathcal S$ is an orbital-space matrix and must not be
confused with the band overlap $O$ or the local integrals $S^x$.
When symmetry is used after the $k$ sum, the same symmetry operation
must be applied to its derivative. Fixed local frames make those
rotation matrices fixed; moving frames add their derivatives.

For fixed weights,
\begin{align}
 \pa U_\kb&=(\pa P_\kb)X+P_\kb\pa X,\label{eq:Uprime}\\
 \pa\mathcal S&=\sum_\kb w_\kb
       [(\pa P_\kb)^\dagger P_\kb+P_\kb^\dagger\pa P_\kb].
 \label{eq:Sprimeorth}
\end{align}
Writing $Y=\mathcal S^{1/2}=X^{-1}$, differentiating
$X\mathcal S X=\Id$ gives the Sylvester equation
\begin{equation}
 Y\pa X+(\pa X)Y=-X(\pa\mathcal S)X.
 \label{eq:Xprime}
\end{equation}
In an eigenbasis of $\mathcal S$ with positive eigenvalues $s_i$,
\begin{equation}
 (\pa X)_{ij}=-\frac{(\pa\mathcal S)_{ij}}
        {\sqrt{s_i}\sqrt{s_j}(\sqrt{s_i}+\sqrt{s_j})}.
\end{equation}
This is well behaved for coincident positive eigenvalues. A singular
global normalization block requires a specified rank treatment;
differentiability can fail when that rank changes.

\subsection{An essential qualification for a fixed-coefficient frame}
The temporary strained states obtained by holding central LAPW
coefficients fixed are not, in general, orthonormal. In that frame the
same physical normalization is
\begin{equation}
 \mathcal S=\sum_\kb w_\kb P_\kb^\dagger O_{W,\kb}^{-1}P_\kb,
 \label{eq:metricS}
\end{equation}
where $O_W$ is the overlap within the chosen window span and $P$ is the
covariant overlap with its basis functions. At the central point $O_W=\Id$,
\begin{equation}
 \pa\mathcal S=\sum_\kb w_\kb\left[
 (\pa P)^\dagger P+P^\dagger\pa P-P^\dagger(\pa O_W)P\right].
 \label{eq:metricSprime}
\end{equation}
Equations~\eqref{eq:Sprimeorth} and \eqref{eq:metricSprime} express
the same normalization in different frames. The last term in
Eq.~\eqref{eq:metricSprime} is \emph{not} an extra term to add to a
correct derivative already evaluated in a moving orthonormal KS basis.

For a pure frame transformation $\psi\mapsto\psi(\Id+L)$,
$\delta P=L^\dagger P$ and $\delta O=L^\dagger+L$ at the reference.
Equation~\eqref{eq:metricSprime} then gives $\delta\mathcal S=0$.
Consequently $\delta X=0$, $\delta U=L^\dagger U$, and the cancellation
in Eq.~\eqref{eq:gaugecancel} holds also for this nonorthogonal
relabeling. This is a useful representation-consistency test, not a
claim that the current code has omitted this term.

\subsection{Projector stress and the direct gradient}
The remaining explicit projector contribution is
\begin{equation}
 \boxed{D_\mu^{\rm proj}
     =2\Rea\sum_\kb w_\kb\Tr_c[C_\kb\pa U_\kb],\qquad
 C_\kb=T\sum_j^{\rm reg}\Sig(z_j)U_\kb^\dagger G_\kb(z_j).}
 \label{eq:projstress}
\end{equation}
There is no dagger between $C$ and $\partial U$: their rectangular
orientations already make the product an orbital matrix. The frequency
sum has the same regularization as the density moments.

One can infer $C$ from $UC=B_L$ using
$C=(U^\dagger U)^+U^\dagger B_L$, where $+$ denotes a
Moore--Penrose pseudoinverse. If $U$ has deficient column rank at one
$\kb$, this reconstructs only the part of $C$ visible through multiplication
by $U$. The Matsubara expression in Eq.~\eqref{eq:projstress} defines
the direct gradient without such an inversion. Agreement of contracted
gradients is a test; it is not a replacement for the derivative of $X$.

\newpart{The projector boundary and other extensions}
\label{sec:extensions}
\subsection{A finite projector at the MT boundary}
Suppose a projector is represented as
$\varphi_\ell(\mathbf r)=\theta(R_a-r_a)\,q_\ell(\mathbf r_a)$,
where $q_\ell$ does not vanish at $R_a$. A moving-domain overlap obeys
the transport formula
\begin{equation}
 \delta\!\int_{B_a}\dd^3r\,\psi^*q
 =\int_{B_a}\dd^3r\,\delta(\psi^*q)
   +\int_{\partial B_a}\dd S\,v_n\psi^*q.
 \label{eq:reynolds}
\end{equation}
Here $B_a$ is the integration ball and $v_n$ is the boundary's normal
displacement per unit perturbation in the coordinates being used.
Equivalently, at fixed coordinate $\mathbf r$,
\begin{equation}
 \delta\theta(R_a-r_a)=\delta(R_a-r_a)
           [\delta R_a+\widehat{\mathbf r}_a\cdot\delta\tau_a].
 \label{eq:theta}
\end{equation}
For fixed-radius physical spheres, $\delta R_a=0$ but the sphere
translates with its center. In atom-centered physical coordinates that
translation is transferred into the derivative of the wave function;
there is no additional translating-boundary term to add a second time.
In a material pullback the same fixed physical sphere becomes a
deformed integration region, so its shape derivative must be handled.

The fact that $q(R_a)\ne0$ means a boundary contribution cannot be
discarded merely by declaring the projector localized. It does not
prove a missing term in an existing radial/matching implementation:
the same contribution may already be contained in its coordinate
transformation. A direct overlap derivative in a specified coordinate
system is the decisive audit. This projector boundary issue is distinct
from basis continuity used in deriving Eq.~\eqref{eq:Hprime}.

\subsection{Local orbitals}
Local orbitals can be included by enlarging the radial primitive set and
the matrices $A,D,S,H,V$. For an extra primitive $u_l(r,E_{\rm LO})$,
$H_{\rm sph}u_l=E_{\rm LO}u_l$. For a linear combination defining an LO,
apply the Hamiltonian to every primitive and differentiate every
coefficient and normalization constraint. Do not reuse the elementary
$\delta_{\lambda,1}$ term for an arbitrary LO without deriving its
radial action. APW+lo also requires its own boundary bookkeeping.

\subsection{Scalar relativity and local-axis constraints}
The cancellation of Jacobian, basis representation, and projector
normalization is representation algebra and survives RELA.
Equations~\eqref{eq:radialeq} and \eqref{eq:T}, however, use the
nonrelativistic differential operator. A RELA extension requires the
actual scalar-relativistic radial Hamiltonian, small-component metric,
and their strain derivatives. Successful RELA fixed-input matrix tests
do not by themselves establish every term in the physical strain path.

If the self-energy is restricted to be diagonal in local correlated
axes, changing those axes changes the approximation. A co-moving
frame requires the derivative of its orbital rotation in $P,U$ and
the corresponding treatment of the imposed self-energy constraint.
The present formulas use fixed local axes. Merely rotating a diagonal
self-energy after truncation is not equivalent to a fully rotationally
covariant calculation.

\newpart{What comes directly from the paper, and what is new here}
\label{sec:paper}
\begin{longtable}{p{0.18\linewidth}p{0.73\linewidth}}
\toprule
Paper location & Role in this derivation\\
\midrule
Eq.\ (10) & DFT Rayleigh quotient. Its normalization cancels the volume
 factor before Eq.\ (12); it cannot be reused unchanged with a general
 interacting energy moment.\\
Eq.\ (11) & Explicit nonrelativistic operator derivative, including
 the kinetic expression in Eq.~\eqref{eq:T}.\\
Eq.\ (13) & Product-rule decomposition. Rewritten exactly as
 $b-\gamma\psi/2+\eta$ in Eq.~\eqref{eq:13}.\\
Eq.\ (20) & DFT incomplete-basis correction. Generalized to independent
 $\rho$ and $\rho_E$ in Eq.~\eqref{eq:20dmft}.\\
Appendix D & Volume-compensated interstitial invariance and local
 radial/angular differentiation.\\
Eq.\ (G1) & Diagonal-occupation expansion of the kinetic bilinear form.
 Its off-diagonal definition is Eq.~\eqref{eq:T}.\\
Eq.\ (G2) & Six diagonal-occupation local contractions. Their matrix
 versions are Eqs.~\eqref{eq:g21}--\eqref{eq:g26}.\\
Eqs.\ (G3--G4) & Matching coefficient derivative $D$, including the
 compensating $\sqrt\Omega$ factor.\\
\bottomrule
\end{longtable}

Blaha's paper does not state the interacting $\rho_E$, $M_\Sigma$,
normalized-projector gradient, or the matrix formulas in the notation
used here. Those are the generalizations derived in this note. The
diagonal DFT result is a necessary limit; it does not by itself determine
all off-diagonal matrix elements.

\subsection{Two places where the integral definitions are safer}
In the supplied PDF, the definition of the radial quantity $W$ below
G4 prints radial derivatives on both radial functions. The direct
derivation of G2 line 5, Eq.~\eqref{eq:Vr}, instead gives
\begin{equation}
 W_{i,L'',j}=\int_0^{R_a}r^3\dd r\,u_i'^*(r)V_{L''}(r)u_j(r),
 \label{eq:W}
\end{equation}
with a derivative only on the bra. This follows because only the bra
wave function is differentiated in the one-sided term; its conjugate
is supplied by the outer $2\Rea$. The printed two-prime definition
appears to be a typographical error. The inspected
\pathf{src/dftKS/atpar.f90} already implements and documents the
single-prime form, so this observation is not a newly discovered code bug.

The third line of G2 also uses compact same-channel radial norm notation.
When its angular operator connects different $l$, use the cross-$l$
overlap of Eq.~\eqref{eq:Sa}, rather than extending same-$l$
orthogonality identities to that case. The inspected stress-integral
code already stores such cross-$l$ integrals. This is an audit of the
meaning of the printed simplification, not evidence that those
integrals are absent from the implementation.

\subsection{Beyond simply changing one band index}
The radial and angular operator structures in G2 need no new
DMFT-specific angular algebra. The essential changes are (i) full
bilinear band matrices, (ii) the independent energy moment, (iii)
consistent volume and band-frame cancellation, and (iv) the derivative
of the normalized correlated projector. Additional radial response,
boundary handling, and relativity are separate approximations to audit.

\newpart{Implementation map and the next audit}
\label{sec:audit}
The paths below are relative to
\pathf{/Users/haule/projects/eDMFT_new/} and refer to the inspected
working tree on 6 November 2026. They locate relevant code; they do
not certify that its total assembly is correct.

\begin{longtable}{>{\raggedright\arraybackslash}p{0.36\linewidth}p{0.56\linewidth}}
\toprule
Source location & Mathematical responsibility / audit question\\
\midrule
\pathf{src/dmft2/cmp_dmft_stress.f90}
 & Contracts density and energy-moment information with the existing
 stress kernels. Check the complete off-diagonal bilinear form against
 $T^\mu,B^\mu,O^{\mu,\eta}$, not only its DFT diagonal.\\
\pathf{cmp_dmft_paired_matching_matrix}
 in the same file & Diagnostic fixed-input matching derivatives,
 with separate spherical, surface, nonspherical, and interstitial
 pieces. These form a distinct derivative convention that must be
 mapped to the material-coordinate formula before substitution.\\
\pathf{src/dmft2/alm_blm_clm.f90},
 \pathf{cmp_alm_blm_clm_strain}
 & Matching derivatives include physical volume-normalization terms.
 Determine explicitly whether each consumer needs $s^\mu$ or Blaha's
 $D^\mu=s^\mu+\gamma_\mu A/2$.\\
\pathf{src/dmft2/stress_integrals.f90}
 & Radial derivative and overlap integrals, including cross-$l$
 $J$ and $I$ channels. Compare with Eqs.~\eqref{eq:Sr} and
 \eqref{eq:Sa}, including small-component conventions.\\
\pathf{src/dftKS/atpar.f90}
 & Nonspherical $V$ and single-prime $W$ radial integrals;
 confirm all bra/ket and LO channels.\\
\pathf{src/dmft2/l2main.f90}
 & Construction and assembly of diagnostic density/moment contractions.
 Keep complete total-stress assembly separate from a successful
 fixed-input derivative check.\\
\pathf{src/dmft2/stress_modules.f90}
 & Diagnostic switches, including
 \pathf{EDMFT_STRESS_PAIRED_DIAG} and
 \pathf{EDMFT_STRESS_PAIRED_ASSEMBLY}. A diagnostic reconstruction
 is not a validated production replacement.\\
\bottomrule
\end{longtable}

\subsection{A concrete order for re-auditing}
\begin{enumerate}[itemsep=5pt]
\item \textbf{Fix the derivative convention in writing.} For every
 routine record its coordinates, physical sphere motion, normalization,
 $c$ treatment, potential held fixed, and Hartree/Ry and shear conventions.
\item \textbf{Construct local integral reference matrices.} Compare all
 six G2 channels with Eqs.~\eqref{eq:g21}--\eqref{eq:g26} before
 assembling the final stress. Test both real and imaginary off-diagonal
 Hermitian density matrices; real test vectors alone cannot expose every
 complex off-diagonal error.
\item \textbf{Recover full $H'$ and $O'$ in one frame.} Check
 Eq.~\eqref{eq:matrixderivs}, including whole-cell volume cancellation,
 MT surfaces, interstitial terms, and all LO contributions. Do not infer
 $O'$ from an energy-linear kernel without establishing its full
 bilinear identity.
\item \textbf{Close the projector derivative in the same frame.}
 Include $X'$, window response, and the metric form of normalization
 when appropriate. Verify $UC=B_L$ and compare the direct gradient
 against pseudoinverse reconstruction. Check projector boundary motion
 through a direct overlap derivative.
\item \textbf{Verify exact invariances.} A pure band-frame rotation must
 cancel between Eqs.~\eqref{eq:bterm} and \eqref{eq:gaugecancel}; a
 fixed-coefficient test has no coefficient-response term. In the
 noninteracting limit $M_\Sigma=0$ and the DFT result must return.
\item \textbf{Only then assemble the physical functional derivative.}
 Keep the KS field-response cancellation and density-functional terms
 consistent with the chosen path. Compare NREL first, then RELA, with
 the prescribed frozen-$\Sigma$ reference functional and relaxed charge.
 Use at least two strain magnitudes and test a shear or nonhydrostatic
 channel as well as isotropic strain.
\end{enumerate}

\key{The present evidence supports deriving and testing the complete
off-diagonal bilinear forms. It does not support adding an isolated
$M_\Sigma O'$ contribution to a total stress whose overlap and projector
bookkeeping has not yet been shown to use the same convention.}

\newpart{What the recent fixed-input tests establish}
\label{sec:tests}
The recent diagnostic contracts the central interacting moments with
Hamiltonian and overlap derivatives formed on a specified fixed-input
path. Define
\begin{equation}
 \rho_o=\rho-\diag\rho,\qquad
 D_\mu^{\rm extra}=\Rea\Tr_b\left[
 \rho_o H'_\mu-
 \left(\tfrac12\{\Es,\rho_o\}+M_\Sigma\right)O'_\mu\right].
 \label{eq:extra}
\end{equation}
The $k$ sum and weights are implicit in this equation only.
The quantity called ``extra'' is a diagnostic difference from a chosen
diagonal reference. It is not the total DMFT stress and not an
independently justified addition to the production result.
The recorded paired diagnostic uses the raw valence-stress assembly
sign, opposite to the energy-derivative sign of this table; for example,
the MnO tag \pathf{PAIRED_MAIN_CORRECTION_COMPLETE_MEAN_GPA} is
$-101.143097512$, corresponding to $+101.143097512$ here.

For isotropic linear strain $h=0.001$, the recorded hydrostatic means
in the physical energy-derivative sign convention are:
\begin{center}
\begin{tabular}{lrrr}
\toprule
Material & Analytic & Central difference & Analytic $-$ FD\\
 & \multicolumn{3}{c}{GPa, $D^{\rm extra}/(3\Omega)$}\\
\midrule
SrVO$_3$ & $-5.701366$ & $-5.695318$ & $-0.006048$\\
FeSe, RELA & $17.599107$ & $17.600406$ & $-0.001299$\\
MnO & $101.143098$ & $101.101274$ & $0.041823$\\
\bottomrule
\end{tabular}
\end{center}
The decomposition of the finite-difference contractions was
\begin{center}
\begin{tabular}{lrr}
\toprule
Material & $\rho_o H'-\tfrac12\{\Es,\rho_o\}O'$ & $-M_\Sigma O'$\\
 & \multicolumn{2}{c}{GPa after $k$ sum and division by $3\Omega$}\\
\midrule
SrVO$_3$ & $-8.066304$ & $2.370986$\\
FeSe, RELA & $-1.690087$ & $19.290493$\\
MnO & $-11.585019$ & $112.686293$\\
\bottomrule
\end{tabular}
\end{center}

These values were obtained before preparation of this PDF; no electronic
structure calculation was rerun for this documentation task.
The large MnO and FeSe components demonstrate why an isolated moment
term is not a reliable physical correction. A complete derivative may
contain large canceling terms.

\textbf{Established within these tests:} the reported analytic
fixed-input contraction is close to the corresponding central matrix
difference in these hydrostatic examples.

\textbf{Not established:} convergence with strain step, all individual
matrix elements, shear channels, the full normalized-projector
derivative in the same frame, or agreement of the complete stationary
free-energy derivative. One-sided derivatives show appreciable curvature
in some cases. The MnO diagnostic process also reported IEEE
invalid/divide-by-zero flags at exit; its printed contractions were
finite, but the flags were not traced. This remains a diagnostic
limitation rather than a clean numerical-validation claim.

\subsection{Recorded locations for reproducibility}
The checker is
\pathf{src/python/check_oneparticle_stress_fd.py}.
The following paths identify the diagnostic snapshots used; temporary
directories are not a durable archive and may later be removed.
\begin{itemize}[itemsep=4pt]
\item SrVO$_3$: center and moments under
 \pathf{/private/tmp/sr_fixedpot_center_case/fd/}; displaced matrices under
 \pathf{/private/tmp/sr_fixed_joint_minus.LMwYDZ/fd/} and
 \pathf{/private/tmp/sr_fixed_joint_plus.J6UBtd/fd/}.
 Reference volume $382.70904\,a_0^3$.
\item FeSe RELA: center and moments under
 \pathf{/private/tmp/fese_rela_fullq.QNEvK6/center/fd/}; displaced
 matrices under
 \pathf{/private/tmp/edmft_fixedpot_oracle_20261106/fese_minus/fd/}
 and the corresponding \pathf{fese_plus/fd/} directory.
 Reference volume $528.96251\,a_0^3$.
\item MnO: center and moments under
 \pathf{/private/tmp/mno_projector_closure.hEJUBK/center/fd/}; displaced
 matrices under
 \pathf{/private/tmp/edmft_fixedpot_oracle_20261106/mno_minus/fd/}
 and the corresponding \pathf{mno_plus/fd/} directory.
 Reference volume $148.11716\,a_0^3$.
 Analytic diagnostic output:
 \pathf{/private/tmp/mno_paired_review.y3ZYV4/MnO.outputdmf2}.
\end{itemize}

\newpart{Compact formula sheet for the next discussion}
\label{sec:sheet}
With all quantities in one representation and the same regularization,
\begin{align}
 D_\mu^{\rm 1p}&=\sum_\kb w_\kb\Rea\Tr_b
                   (\rho\pa H-\rE\pa O),\\
 \rE&=\tfrac12\{\Es,\rho\}+M_\Sigma,\\
 \pa\psi&=b-\tfrac12\gamma_\mu\psi+\eta,\\
 \gamma_\mu\Tr(\rho H-\rE O)
     +\mathcal P_\mu[-\gamma_\mu\psi/2]&=0,\\
 \eta_n^\mu&=\sum_i[D^\mu_{in}u_iY_i
               +A_{in}r u_i'p_\mu Y_i
               +A_{in}u_i\mathscr A_\mu Y_i],\\
 D^\mu&=s^\mu+\tfrac12\gamma_\mu A,\\
 \sigma_\mu^{\rm corr}
 &=\frac2\Omega\Rea\sum_{a\kb}w_\kb
       \sum_{x=0,r,\mathrm{ang}}
       \Tr_b[\rho Q_x^\dagger(H^x+V^x)A
                        -\rE Q_x^\dagger S^x A],\\
 Q_0&=D^\mu,\qquad Q_r=Q_{\mathrm{ang}}=A,\\
 D_\mu^{\rm proj}&=2\Rea\sum_\kb w_\kb\Tr_c[C\pa U],\\
 U&=PX,\qquad \pa U=(\pa P)X+P\pa X,\\
 Y\pa X+(\pa X)Y&=-X(\pa\mathcal S)X,\qquad Y=\mathcal S^{1/2},\\
 UC&=\rE-\Es\rho=B_L,\\
 \mathcal P_\mu[b]+D_\mu^{\rm proj,frame}&=0
       \quad\text{for a consistent pure frame change}.
\end{align}
In the compact $Q_x$ notation, $H^0,V^0,S^0$ pair with $D^\mu$;
the other two channels pair with $A$. This sheet is a restatement of
the derived formulas, not an independent shortcut for the omitted
potential, core, ion, boundary, or relativistic terms.

\begin{thebibliography}{9}
\bibitem{blaha}
Belbase, Tr\"oster, and Blaha,
\emph{Stress tensor in the linearized augmented plane wave method},
Physical Review B \textbf{104}, 174113 (2021).
\href{https://doi.org/10.1103/PhysRevB.104.174113}{doi:10.1103/PhysRevB.104.174113}.
Source inspected: \pathf{/Users/haule/papers_to_read/Blahas_stress_tensor.pdf}.
Equations (10)--(13), (19)--(20), Appendix D, and Appendix G.

\bibitem{forces}
K. Haule, project documentation on the DFT+eDMFT force functional,
\pathf{/Users/haule/projects/eDMFT/src/doc/Forces_LAPW.pdf}.
Background for the stationary-functional and normalized-projector
discussion; the strain derivation is given explicitly in this note.

\bibitem{previous}
Earlier project working note,
\pathf{src/doc/stress/edmft_stress_implemented_equations.tex}
and \pathf{Analytic_Stress_DFT_eDMFT_Implemented_Equations.pdf},
5 November 2026. The present note supplements that record and makes
the Eq.\ (13) cancellation and local G2 matrix definitions explicit.

\bibitem{source}
Local working tree \pathf{/Users/haule/projects/eDMFT_new/},
read-only source inspection, 6 November 2026.
The source map in Sec.~\ref{sec:audit} identifies relevant routines;
implementation correctness remains subject to the stated audit.
\end{thebibliography}
\end{document}
